% GENERATED FILE. Edit canonical lessons, then run npm run generate:books.
% canonical-source-hash: fnv1a64:7369ee76a1251d69
% canonical-lessons: GU-C522-khayni, GU-C522-khadvu, GU-C522-chinvu, GU-C522-palaalvu, GU-C522-dalvu

\chapter{Pound, Grate and Soak}
\label{ch:gu-pound-grate-and-soak}
\hlchaptermodality{522}

\begin{chapteropening}
\textbf{By the end of this chapter:} \emph{I can say how food is pounded, grated, soaked and ground.}

Complete the last lesson of chapter 522: I can say how food is pounded, grated, soaked and ground.
\end{chapteropening}

\section[\texorpdfstring{\gu{ખાયણી} (khāyṇī)}{ખાયણી (khāyṇī)}]{\gu{ખાયણી} (khāyṇī) — a mortar (for pounding)}

\label{lesson:GU-C522-khayni}



\begin{quote}
Before the new one: say the Gujarati for to boil over, then the Gujarati for to whisk.
\end{quote}

\subsection*{You'll want to know: \gu{ખાયણી}}
\textbf{\gu{ખાયણી}} — \emph{khāyṇī} — \textquotedblleft{}a mortar (for pounding)\textquotedblright{}.

\begin{quote}
\gu{દાદીની} \gu{ખાયણી} \gu{પથ્થરની} \gu{છે}.
\end{quote}

\begin{quote}
\emph{dādīnī khāyṇī paththarnī che.}
\end{quote}

\begin{quote}
Grandmother's mortar is made of stone.
\end{quote}

A feminine noun.

\begin{cousinweb}
\textbf{\gu{ખાયણી}} (\emph{khāyṇī}) is filed by CDIAL under Sanskrit \emph{khaṇḍana}, \textquotedblleft{}breaking\textquotedblright{}, with the Gujarati word for a mortar (CDIAL 3796).
\end{cousinweb}

\begin{grammarlens}[title={What you've built}]
One more word for describing what you lived and felt.
\end{grammarlens}

\subsection*{Your turn}
\begin{itemize}
\raggedright
  \item \emph{Say it:} \emph{khāyṇī}
  \item \emph{Say it:} \emph{khāyṇī}, once more
  \item \emph{Say it:} \emph{dādīnī khāyṇī paththarnī che.} — Grandmother's mortar is made of stone.
  \item \emph{Recall:} say the Gujarati for to boil over, then the Gujarati for to whisk, then say \emph{khāyṇī} again
\end{itemize}

\begin{tcolorbox}[breakable,colback=teal!4,colframe=teal!35!black,title={Before you move on}]
What does \gu{ખાયણી} mean? (\textquotedblleft{}A mortar (for pounding)\textquotedblright{}.) Say it once more.
\end{tcolorbox}
\section[\texorpdfstring{\gu{ખાંડવું} (khā̃ḍvũ)}{ખાંડવું (khā̃ḍvũ)}]{\gu{ખાંડવું} (khā̃ḍvũ) — to pound}

\label{lesson:GU-C522-khadvu}



\begin{quote}
Before the new one: say the Gujarati for to whisk, then the Gujarati for a mortar.
\end{quote}

\subsection*{You'll want to know: \gu{ખાંડવું}}
\textbf{\gu{ખાંડવું}} — \emph{khā̃ḍvũ} — \textquotedblleft{}to pound\textquotedblright{}.

\begin{quote}
\gu{દાદીએ} \gu{ખાયણીમાં} \gu{આદુ} \gu{ખાંડ્યું}.
\end{quote}

\begin{quote}
\emph{dādīe khāyṇīmā̃ ādu khā̃ḍyũ.}
\end{quote}

\begin{quote}
Grandmother pounded the ginger in the mortar.
\end{quote}

Past, with \gu{તેણે} (he or she …): \gu{તેણે} \gu{આદુ} \gu{ખાંડ્યું}, \emph{teṇe ādu khā̃ḍyũ}.

\begin{cousinweb}
\textbf{\gu{ખાંડવું}} (\emph{khā̃ḍvũ}) is inherited from Sanskrit \emph{kháṇḍatē}, \textquotedblleft{}breaks\textquotedblright{}, and CDIAL lists the Gujarati verb, \textquotedblleft{}to break, pound\textquotedblright{} (CDIAL 3795). Sugar, \textbf{\gu{ખાંડ}}, comes from the same root: sugar in pieces.
\end{cousinweb}

\begin{grammarlens}[title={What you've built}]
One more word for describing what you lived and felt.
\end{grammarlens}

\subsection*{Your turn}
\begin{itemize}
\raggedright
  \item \emph{Say it:} \emph{khā̃ḍvũ}
  \item \emph{Say it:} \emph{khā̃ḍvũ}, once more
  \item \emph{Say it:} \emph{dādīe khāyṇīmā̃ ādu khā̃ḍyũ.} — Grandmother pounded the ginger in the mortar.
  \item \emph{Recall:} say the Gujarati for to whisk, then the Gujarati for a mortar, then say \emph{khā̃ḍvũ} again
\end{itemize}

\begin{tcolorbox}[breakable,colback=teal!4,colframe=teal!35!black,title={Before you move on}]
What does \gu{ખાંડવું} mean? (\textquotedblleft{}To pound\textquotedblright{}.) Say it once more.
\end{tcolorbox}
\section[\texorpdfstring{\gu{છીણવું} (chīṇvũ)}{છીણવું (chīṇvũ)}]{\gu{છીણવું} (chīṇvũ) — to grate}

\label{lesson:GU-C522-chinvu}



\begin{quote}
Before the new one: say the Gujarati for a mortar, then the Gujarati for to pound.
\end{quote}

\subsection*{You'll want to know: \gu{છીણવું}}
\textbf{\gu{છીણવું}} — \emph{chīṇvũ} — \textquotedblleft{}to grate\textquotedblright{}.

\begin{quote}
\gu{મીરાએ} \gu{ગાજર} \gu{છીણ્યું}.
\end{quote}

\begin{quote}
\emph{Mīrāe gājar chīṇyũ.}
\end{quote}

\begin{quote}
Mira grated the carrot.
\end{quote}

Past, with \gu{તેણે} (he or she …): \gu{તેણે} \gu{ગાજર} \gu{છીણ્યું}, \emph{teṇe gājar chīṇyũ}.

\begin{cousinweb}
\textbf{\gu{છીણવું}} (\emph{chīṇvũ}) is filed by CDIAL under Sanskrit \emph{chinátti}, \textquotedblleft{}cuts off\textquotedblright{}, with a Gujarati \emph{chiṇvũ}, \textquotedblleft{}to slit, peel\textquotedblright{} (CDIAL 5045); grating is probably the same cutting, into fine shreds.
\end{cousinweb}

\begin{grammarlens}[title={What you've built}]
One more word for describing what you lived and felt.
\end{grammarlens}

\subsection*{Your turn}
\begin{itemize}
\raggedright
  \item \emph{Say it:} \emph{chīṇvũ}
  \item \emph{Say it:} \emph{chīṇvũ}, once more
  \item \emph{Say it:} \emph{Mīrāe gājar chīṇyũ.} — Mira grated the carrot.
  \item \emph{Recall:} say the Gujarati for a mortar, then the Gujarati for to pound, then say \emph{chīṇvũ} again
\end{itemize}

\begin{tcolorbox}[breakable,colback=teal!4,colframe=teal!35!black,title={Before you move on}]
What does \gu{છીણવું} mean? (\textquotedblleft{}To grate\textquotedblright{}.) Say it once more.
\end{tcolorbox}
\section[\texorpdfstring{\gu{પલાળવું} (palāḷvũ)}{પલાળવું (palāḷvũ)}]{\gu{પલાળવું} (palāḷvũ) — to soak}

\label{lesson:GU-C522-palaalvu}



\begin{quote}
Before the new one: say the Gujarati for to pound, then the Gujarati for to grate.
\end{quote}

\subsection*{You'll want to know: \gu{પલાળવું}}
\textbf{\gu{પલાળવું}} — \emph{palāḷvũ} — \textquotedblleft{}to soak\textquotedblright{}.

\begin{quote}
\gu{રાત્રે} \gu{મગ} \gu{પાણીમાં} \gu{પલાળો}.
\end{quote}

\begin{quote}
\emph{rātre mag pāṇīmā̃ palāḷo.}
\end{quote}

\begin{quote}
Soak the mung beans in water overnight.
\end{quote}

Past, with \gu{તેણે} (he or she …): \gu{તેણે} \gu{મગ} \gu{પલાળ્યા}, \emph{teṇe mag palāḷyā}.

\begin{cousinweb}
\textbf{\gu{પલાળવું}} (\emph{palāḷvũ}) is the causative partner of \textbf{\gu{પલળવું}}, \textquotedblleft{}to get soaked\textquotedblright{}, and CDIAL lists it under \emph{*pralālayatē} as \textquotedblleft{}to make wet, soak\textquotedblright{} (CDIAL 8759).
\end{cousinweb}

\begin{grammarlens}[title={What you've built}]
One more word for describing what you lived and felt.
\end{grammarlens}

\subsection*{Your turn}
\begin{itemize}
\raggedright
  \item \emph{Say it:} \emph{palāḷvũ}
  \item \emph{Say it:} \emph{palāḷvũ}, once more
  \item \emph{Say it:} \emph{rātre mag pāṇīmā̃ palāḷo.} — Soak the mung beans in water overnight.
  \item \emph{Recall:} say the Gujarati for to pound, then the Gujarati for to grate, then say \emph{palāḷvũ} again
\end{itemize}

\begin{tcolorbox}[breakable,colback=teal!4,colframe=teal!35!black,title={Before you move on}]
What does \gu{પલાળવું} mean? (\textquotedblleft{}To soak\textquotedblright{}.) Say it once more.
\end{tcolorbox}
\section[\texorpdfstring{\gu{દળવું} (daḷvũ)}{દળવું (daḷvũ)}]{\gu{દળવું} (daḷvũ) — to mill grain}

\label{lesson:GU-C522-dalvu}



\begin{quote}
Before the new one: say the Gujarati for to grate, then the Gujarati for to soak.
\end{quote}

\subsection*{You'll want to know: \gu{દળવું}}
\textbf{\gu{દળવું}} — \emph{daḷvũ} — \textquotedblleft{}to mill grain\textquotedblright{}.

\begin{quote}
\gu{દાદી} \gu{ઘઉં} \gu{દળીને} \gu{લોટ} \gu{બનાવતી} \gu{હતી}.
\end{quote}

\begin{quote}
\emph{dādī ghaũ daḷīne loṭ banāvtī hatī.}
\end{quote}

\begin{quote}
Grandmother used to grind wheat into flour.
\end{quote}

Past, with \gu{તેણે} (he or she …): \gu{તેણે} \gu{ઘઉં} \gu{દળ્યા}, \emph{teṇe ghaũ daḷyā}.

\begin{cousinweb}
\textbf{\gu{દળવું}} (\emph{daḷvũ}) is inherited from Sanskrit \emph{dálati}, \textquotedblleft{}bursts, splits\textquotedblright{}, and CDIAL lists the Gujarati verb, \textquotedblleft{}to crush, grind\textquotedblright{} (CDIAL 6216).
\end{cousinweb}

\begin{grammarlens}[title={What you've built}]
One more word for describing what you lived and felt.
\end{grammarlens}

\subsection*{Your turn}
\begin{itemize}
\raggedright
  \item \emph{Say it:} \emph{daḷvũ}
  \item \emph{Say it:} \emph{daḷvũ}, once more
  \item \emph{Say it:} \emph{dādī ghaũ daḷīne loṭ banāvtī hatī.} — Grandmother used to grind wheat into flour.
  \item \emph{Recall:} say the Gujarati for to grate, then the Gujarati for to soak, then say \emph{daḷvũ} again
\end{itemize}

\begin{tcolorbox}[breakable,colback=teal!4,colframe=teal!35!black,title={Before you move on}]
What does \gu{દળવું} mean? (\textquotedblleft{}To mill grain\textquotedblright{}.) Say it once more.
\end{tcolorbox}
